\section{Related Work}
\label{sec:related}

\nospacestitle{Reuse {\kvcache} across requests via {\kvcache} cache. \,}
Reusing {\kvcache} for accelerating LLM serving
has been widely studied~\cite{vllm,chunkattention,cachedattention,promptcache,sglang,cacheblend} and used in commercial LLM serving~\cite{openaiapi,genimiapi,claudeapi,mooncake}.
Currently, production systems only treat cache hits by prefix matches~\cite{chunkattention,sglang,cachedattention,openaiapi,genimiapi,claudeapi,mooncake}, 
because it preserves the original algorithm of model inference with no accuracy loss.
A number of studies~\cite{promptcache,cacheblend,hu2024epic} have studied 
methods for non-prefix {\kvcache} cache. 
We can revisit our characterization once they have matured 
and been deployed in production.

\stitle{Other {\kvcache}-related optimizations. \,}
Emerging works\cite{streamingllm,h2o,infinigen,pyramidkv} propose runtime {\kvcache}
compression and deletion methods to reduce {\kvcache} size.
StreamingLLM~\cite{streamingllm} keeps
only a finite attention window for recent tokens combined with a few initial tokens.
H2O~\cite{h2o} only involves tokens that contribute most to the attention score for inference.
InfiniGen~\cite{infinigen} speculatively selects tokens critical for attention scores and drops others.
KVQuant~\cite{KVQuant} enables 3-bit KV cache quantization with minimal perplexity degradation 
for large language models, achieving fast million-length context inference.
These methods require fewer {\kvcache} per-request,
albeit with accuracy degradation.
Our study is compatible with these works:
if an uncompressed/full {\kvcache} is reusable,
a compressed or deleted version would also be reusable.
%\TODO{add more about the recent sparse attention works.}

%\stitle{Cache in web-service and serverless platform. }
\stitle{Optimizing caching policies. \,}
Optimizing caching has long been studied in the literature,
in general-purpose caching policies~\cite{lruk,slru,twoq,eelru,lrfu,lirs,arc,mq,car,clockpro,DBLP:journals/tos/EinzigerEFM22,lhd,cacheus,sieve,cherkasova1998improving}
or specific domains~\cite{yang2020twemcache,berg2020cachelib,icebreaker,fasscache}.
We continue this line of research,
and leverage our characterized {\kvcache} reuse properties
for optimizing {\kvcache} cache policies.

\stitle{Optimizing LLM serving. \,}
We continue the line of research in optimizing the performance of LLM serving systems~\cite{DBLP:conf/osdi/ZhongLCHZL0024,DBLP:journals/corr/abs-2404-09526,DBLP:conf/sosp/KwonLZ0ZY0ZS23,alpaserve,DBLP:conf/osdi/YuJKKC22,DBLP:conf/isca/PatelCZSGMB24,cachedattention,DBLP:journals/corr/abs-2412-17246},
with a particular focus on characterizing serving workloads
for {\kvcache} cache.
Our work is orthogonal to optimizations other than {\kvcache} cache.
\iffalse
% Cache replacement algorithm is essential to computer systems, and has been researched
% for decades\cite{lruk,slru,lirs,arc,lhd,cacheus,sieve}.
The design of a cache replacement policy needs to take into account the
characteristics of the workload~\cite{lruk,slru,lirs,arc,lhd,cacheus,sieve}.
Therefore, access traces analysis for real-world production systems can help to understand the workload characteristics~\cite{yang2020twemcache,berg2020cachelib,shahrad2020serverless,fasscache,icebreaker}.
% \cite{DBLP:conf/osdi/YangYR20,DBLP:conf/osdi/BergBMGGLUCBHG20} characterize cache workloads in production web-service systems.
% They show that skewed Zipf distribution is common for popularity distribution, and requests for some entry are very bursty.
Studies of web services~\cite{yang2020twemcache,berg2020cachelib}
have shown that web requests have a Zipf distribution, with bursts of hotspots.
Studies of serverless platforms~\cite{shahrad2020serverless,fasscache,icebreaker} design keep-alive policies for function caching to accelerate function invocations based on the observation that function calling frequencies vary.
% \cite{shahrad2020serverless} provides a thorough characterization of Function as a Service (FaaS) workloads
% in production system, and reveals varieties in function calling frequencies.
% Based on this observation, multiple works\cite{shahrad2020serverless,fasscache,icebreaker} design keep-alive
% policies for function caching to accelerate function invocations.
Our work performs trace analysis for different work scenarios (toB and toC) for LLM inference and optimizes {\kvcache} cache based on workload characteristics.
We find similarities and differences between the workloads in the LLM inference scenario and the previous cache systems. \SSJ{Some of the similarities and differences can be specifically summarised to highlight our contribution.}
\fi
